\chapter{Аудит смешанных мостов морщинки и дислокации}
\label{ch:v10-particle-wrinkle-dislocation-mixed-bridge-audit}

Предыдущий гейт установил, что топологический дефект и конечный
энергетический профиль совместимы, но лежат в прямой сумме. Теперь проверены
одиннадцать возможных механизмов смешивания: унаследованная прямая сумма,
скалярное произведение следов, проекторная кривизна, произведение
тёплицевой и пространственной границ, градуированный моритов коннектор,
хопфово-черново спаривание, профиль Каллиаса, клеточная инцидентность K43,
ковариация ванны, хиггсовская смена ранга и target-loaded полюс.

Шесть критериев требуют общий типизированный носитель, ненулевой смешанный
блок, сохранение индекса, конечную локализацию, полюс того же оператора и
унаследованность без загрузки ответа. Матрица аудита имеет ранг $6$, а
оценки равны
\begin{equation}
 (3,3,4,4,3,4,5,4,4,3,5).
\end{equation}
Строгих и унаследованных проходов нет.

Две особенно естественные ветви закрыты ранее накопленными тождествами.
В градуированном произведении смешанные члены квадрата сокращаются, а
центрированная моритова кривизна аддитивна. Поэтому оба mixed block равны
нулю, как и в reopening-гессиане $\operatorname{diag}(2,0)$.

Лучший нетавтологический кандидат --- профиль Каллиаса. Условно он
сохраняет индекс, локализует проектор ранга $15$ и даёт положительный
невырожденный гессиан
\begin{equation}
 H_{\rm C}=\begin{pmatrix}2&1\\1&1\end{pmatrix},
 \qquad
 \operatorname{Spec}H_{\rm C}
 =\left\{\frac{3-\sqrt5}{2},\frac{3+\sqrt5}{2}\right\}.
\end{equation}
Но его пространственный Clifford-носитель, профиль массы и нормировка не
следуют из текущего конечного оператора. Target-loaded сопоставление также
получает пять баллов, но запрещено именно потому, что вставляет искомый
полюс.

\begin{theorem}[Каллиасов профиль является минимальным полным кандидатом]
Среди проверенных нетавтологических мостов только профиль Каллиаса
одновременно допускает ненулевое смешивание, сохранение индекса, конечную
локализацию и спектральный полюс одного оператора.
\end{theorem}
\begin{nogocriterion}[Условная полнота не равна происхождению]
Пока spinor--Clifford carrier и массовый профиль не вложены в накопленный
носитель, каллиасова конструкция является внешней архитектурой, а не
выведенной частицей.
\end{nogocriterion}
\begin{protocol}\begin{enumerate}
\item покрытие кандидатов: \textbf{$11/11$};
\item строгие проходы: \textbf{$0/11$};
\item унаследованные проходы: \textbf{$0/11$};
\item лучший кандидат: \textbf{профиль Каллиаса, $5/6$};
\item физическое происхождение: \textbf{$0/3$};
\item следующий гейт: \textbf{вложение каллиасова носителя}.
\end{enumerate}\end{protocol}
Машинный сертификат сохранён в
\path{s2t_v10_particle_wrinkle_dislocation_mixed_bridge_candidate_audit_gate_results.json}.